\chapter{Follow prediction, validation and seeking}
\MARSChapterHeading{Follow prediction, validation and seeking}
\label{revision-lab:lossless-video}

Read Chapter 25 with this lab. Import the native lesson \emph{Follow video
history and exclusive bounds}, built from
\path{revision/lessons/034-lossless-video}. Its six experiments lead to one
continuing independent source, \texttt{independent/video-change-bounds.asm}.
Only that bounds function is unfinished. The four helper experiments retain
the supplied codec bodies and relevant static data; their small callers appear
first. Two further sources model deadline arithmetic and selected coordinates.
MLV1 is the local teaching format. The following H.263 work targets an actual
interoperable format.

Answer the predictions inline, then use ordinary Debug at the named boundaries.
Inspect state before executing each anchored instruction and finish the process
before the next activity. If you alter state or backstep while exploring, start
a fresh forward session to record the required observations. A source pointer
identifies a lesson location; it does not establish an executed instruction.
The required experiments read no host file and request no host clock. Hints and
free-writing availability follow the instructor's policy.

The complete standalone example remains in the book's
\texttt{revision/labs/files} directory as \texttt{lossless-video.asm}.
The original construction capstone is in
\path{packages/034-lossless-video/}, with its separate callable contracts. Those materials support additional experiments after
the native lesson; they are not extra steps needed to implement change bounds.

The standalone full-codec file keeps earlier example callers under distinct labels, with the
complete file/player caller named \texttt{main}. Launch that entry for the
end-to-end pipeline. Use the public function names when checking an individual
stage; do not combine the separate historical examples as additional definitions.
If a wrapper places its selected \texttt{main} after another executable example,
export it with \texttt{.globl main} and enable start-at-main. That launch option
uses the exported entry; a grading call can select a callable label directly.
The example's compact arenas have 262144-byte capacities. This accommodates the
supplied clips without reducing the callable format limits. Maximum-size grading
supplies its own separate storage.

\section{Predict one update before applying a record}

Begin with \texttt{worked/xor-alias-tail.asm}. Five native word pairs differ
by hexadecimal \texttt{80000000}, zero, \texttt{10}, \texttt{ffffffff} and
one. At \texttt{deltaKernel\_video\_group\_counted}, inspect all 128 bits of
\texttt{w2} and its four numerical word lanes. Three lanes differ. The high-bit
lane is nonzero even though the same bit pattern is negative when interpreted
as signed. At \texttt{lesson\_xor\_result}, the scalar tail raises the count to
four, and the first guard remains hexadecimal \texttt{a5}. The empty call then
succeeds with deliberately unusable pointers before dereferencing them.

Figure~\ref{fig:revision-lossless-native-xor} shows the current vector boundary
in the actual debugger. Count the nonzero words before including the tail.

\begin{figure}[!ht]
  \centering
  \begin{tikzpicture}[alt={Three crops of one actual light-theme debugger photograph. A identifies the current paused frame at xor-alias-tail.asm line 86. B shows the full written w2 vector. C shows four 32-bit lanes, beginning with lane zero: hexadecimal 80000000, 00000000, 00000010 and ffffffff. Three are nonzero, including the high-bit-only value and all-ones value. This is before the cursor advances to the scalar fifth-word tail.}]
    \node[anchor=north west,inner sep=0,font=\small\bfseries] (title) at (0,0) {Count nonzero words before the scalar tail};
    \node[anchor=north west,inner sep=0] (frame) at ([yshift=-3mm]title.south west) {\MARSRevisionCrop{revision/captures/lossless/lossless-native-xor.png}{0.376\linewidth}{0.03828125}{0.22222222}{0.77812500}{0.75111111}};
    \node[anchor=east,font=\small\bfseries] at ([xshift=-1mm]frame.west) {A};
    \node[anchor=north west,inner sep=0] (vector) at ([yshift=-3mm]frame.south west) {\MARSRevisionCrop{revision/captures/lossless/lossless-native-xor.png}{0.96\linewidth}{0.24843750}{0.09777778}{0.28281250}{0.87777778}};
    \node[anchor=east,font=\small\bfseries] at ([xshift=-1mm]vector.west) {B};
    \node[anchor=north west,inner sep=0] (lanes) at ([yshift=-3mm]vector.south west) {\MARSRevisionCrop{revision/captures/lossless/lossless-native-xor.png}{0.96\linewidth}{0.26171875}{0.03333333}{0.26953125}{0.94333333}};
    \node[anchor=east,font=\small\bfseries] at ([xshift=-1mm]lanes.west) {C};
  \end{tikzpicture}
  \caption{The current frame is the counted vector group at line 86 (A).
    The full \texttt{w2} value (B) and word lanes (C) come from the same photograph.
    Read lane zero first: \texttt{80000000} is nonzero, zero is equal,
    \texttt{00000010} is nonzero, and \texttt{ffffffff} is nonzero. Thus three
    of these four words differ. A signed-positive test would miss two of them.
    The vector display places its low lane on the right; the lane list begins
    with that low lane. The fifth input word remains for the scalar tail.}
  \label{fig:revision-lossless-native-xor}
\end{figure}


The result exactly aliases the left input. Both input vectors are loaded before
the vector store, retaining the original values needed for that group. This
explains the admitted exact alias; it does not prove arbitrary partial overlap
safe. The complete word comparison is also distinct from serialized RGB
validation: a numerical kernel can compare every bit without declaring every
word a valid file-format pixel.

Next use \texttt{worked/record-publication.asm}. Its first delta changes the
initial numerical pixel from hexadecimal \texttt{000061c2} to
\texttt{000161c2}. At \texttt{lesson\_record\_committed}, the context expects
sequence two and publishes kind one and 16384 image bytes. The next candidate
also has sequence two and a recomputed CRC, but its final zero/zero command
cannot advance. Compare it with the context at
\texttt{lesson\_record\_before\_retry}. After rejection, status is one and
count zero; the preceding successful pixel and context fields remain. The
supplied \texttt{worked/late-record.bin} contains those exact candidate bytes.
One retained pixel does not prove that every image byte was preserved.

For the additional artifact experiment, generate the moving fixture in a fresh
directory. Its observations describe the
records, while its PPM files expose independently planned RGB frames. Identify
which stored values are big-endian file fields and which in-memory values are
native words. Explain why four native bytes per pixel become three serialized
bytes and why XOR compares numerical words independently of memory byte order.

Trace the eleven-byte first delta. After skip zero/change one, write the pixel
cursor and the first updated word. Then consume skip 4095/change zero. The
cursor ends at 4096, and no terminator is needed. For the two-by-two patch, the
linear indices are 325, 326, 389 and 390. Explain why the second skip is 62,
using the cursor after the preceding two changed pixels rather than their start.

Now make the last command zero/zero and repair the CRC. A checksum failure
would not test the command walk. Predict the status and every image/context
byte that must remain unchanged. Inspect \texttt{record\_call\_return} in the
standalone record example, or invoke \texttt{video\_apply\_record} through its
public checks. Locate the first pixel store and establish why no later format
validation can still fail after it.

Extend the XOR experiment by comparing a right-input alias and a disjoint
result with its supplied left-input alias. Explain why both source words must be loaded before their
store. Include a zero-count call with unusable pointers and a difference of
hexadecimal \texttt{80000000}. Neither a positive signed comparison nor an
extra alignment condition belongs in the count contract.

\section{Make the cost model include its headers}

Use \texttt{worked/encoding-cost.asm}. Its four changed pixels occupy three
runs: pixel zero, pixels 325 and 326, and pixel 4095. With $k=3$ and $m=4$,
$4k+3m=24$ payload bytes; the record header makes forty. Compare the payload
field and complete returned length at \texttt{lesson\_encoded}. A successful
encode commits current pixels as previous. Reusing that current image at the
next nonperiodic sequence therefore needs only a four-byte unchanged command
and a sixteen-byte record header. At \texttt{lesson\_encoded\_reuse}, the
complete length is twenty and the next sequence is three.

Inspect \texttt{lesson\_encoding\_rejected} after the call.
The caller's capacity is nineteen, below the minimum twenty-byte record.
The previous successful length remains twenty and sequence remains three. The capacity stays nineteen:
it is an explicit input edit made by the caller, not a value that this failure
must undo.

Now derive $4k+3m$ for $k$ commands and $m$ changed pixels. For an unchanged frame,
a tail-only change and a first-pixel-only change, derive payload counts four,
seven and eleven. Each record also needs sixteen header bytes. Then compare
an all-changed frame with an alternating change pattern. Counting changed words
alone cannot decide which representation is smaller.

The worked encoder validates all source colors before writing. Change the high
byte of the final source word and check that output, previous image and final
context remain unchanged. For a nonperiodic record, also validate previous-image
colors; at a mandatory full position, old pixels are irrelevant. A scratch XOR
array may change on failure, but the next committed sequence must not.

For a selected-result capacity experiment, restart the first-call state and
set capacity to 39, one below its forty-byte result. Retry with enough space.
After successful reuse, capacity 39 would be ample for the twenty-byte result;
changing the same number without restoring the same history tests a different
case.
The retry must compare against the original previously committed image. Do not
initialize an independent decoder with the encoder's newly updated previous
image; preserve the actual old image and decode the produced record against it.
Otherwise a broken transition can be hidden by giving the decoder its expected
answer as starting state.

Additional full records are legal. Public feedback accepts their structure and
reconstructed pixels instead of requiring the worked greedy byte sequence.
If a course asks specifically for compression gains or vector code, state that
extra objective separately from format correctness.

\section{Keep a directory tied to its validated file}

The native source is \path{worked/directory-seek.asm}. Its embedded
moving file requires no path edit.
The \texttt{lesson\_directory\_ready} stop shows eight native offsets
occupying 64 bytes. Key interval is four and frame interval is forty milliseconds. The
native caller and the standalone file example use separate buffers, so follow
the source selected by the activity rather than assuming one address for both.

Calculate all eight record offsets from the moving fixture's payload lengths.
Add the sixteen-byte record header each time. The last end must be 24795.
Directory entries are native offsets, not addresses or pixel indices. Before
adding one to a file base, establish that its record header and exact interval
belong to the retained file.

Index the whole file, then alter a late record in a separate candidate. Failure
must preserve the previous successful directory and context. Scratch may have
changed. Explain how the two-pass publisher obtains this guarantee without
requiring a second complete file-sized allocation. Keep the admitted source
immutable throughout both passes and for the directory's lifetime.

At \texttt{lesson\_sought}, the native caller has sought to six from an
unrelated display image. Its returned image extent is 16384 bytes. Inspect the
published target six, replay start four and applied-record count three. Moving data must replay records
four, five and six. The extra-full fixture stores five independently, so it can
start there and replay only five and six. Both produce the same requested words.
Repeat with targets zero and three. Starting from a periodic boundary alone
misses optional full frames; applying only the target delta misses its history.

The native caller next requests target eight, outside indices zero through
seven. At \texttt{lesson\_seek\_rejected}, status is one and count zero,
while the last successful result fields remain six, four and three. Those fields
alone do not establish complete image preservation.

A seek reconstructs in separate scratch and copies only after success. Reject
an out-of-range target, an offset before byte 24, a short record interval and
a misaligned native directory before the corresponding invalid access. The
whole-file validation precondition still applies: local interval checks do not
turn arbitrary directory bytes into a trusted index.

\section{Check the schedule without measuring the host by accident}

Use \texttt{worked/anchored-deadlines.asm} for the native arithmetic model.
At \texttt{lesson\_early\_wake}, deadline 1040 gives waits thirty and five
for readings 1010 and 1035. At \texttt{lesson\_late\_reading}, reading
1047 gives lateness seven and next deadline 1080. The final observation is
1320, or eight forty-millisecond intervals after anchor 1000. This source makes
no host-clock request; it illustrates the relation between readings and deadlines.
The separate scripted player checks and optional real playback follow below.

For anchor 1000 and forty-millisecond intervals, write deadlines through the
last frame. An unchanged frame and the final frame each require their interval.
For observed readings 1010, 1035 and 1040, predict sleeps thirty and five.
For a late reading 1047, record lateness seven and keep the next deadline 1080.
Do not move the anchor forward by seven or silently discard a frame.

The scripted grading clock supplies exact observations and records sleep
requests without real waiting. It can check backward readings, high-bit time
halves, deadline overflow and early wakeups reproducibly. Real service thirty
and thirty-two calls can be late. Record actual lateness separately from
logical duration and from simulator instruction counts. A screen photograph
cannot prove the instant at which the host presented a frame.

Pause before a clock request after an accepted frame and inspect the complete
working image. On a later playback failure, the accepted count and last visible
image remain meaningful, while final context outputs are unchanged. Compare
this with seek failure, which publishes no new image. Avoid a single blanket
claim about output rollback across all interfaces.

The store routine follows positive returned byte counts and closes its acquired
descriptor on every completed path. It does not promise atomic replacement.
Its published transport contract includes a readable guard byte after the source
extent because the current educational write service prefetches one byte; that
guard is not transferred. Keep this service boundary separate from the exact
selected-word windows of your own numerical kernels. Reloading establishes bytes;
index the reloaded content before playing it.

\section{Implement minimal exclusive change bounds}

Use \texttt{worked/exclusive-coordinates.asm} before implementation. Division
and remainder map indices one and four in a two-column image to $(1,0)$ and
$(0,2)$. The containing exclusive ends are two and three, giving area six;
the changed count is two. This source considers only selected positions. Your
function must find extrema across every selected input word. Its final pause
compares full native width hexadecimal \texttt{0000000100000001} with its low
word one. The complete width exceeds 64 and must be rejected.

Replace the deliberate \texttt{break} at \texttt{video\_change\_bounds} in the
continuing \texttt{independent/video-change-bounds.asm}. Arguments are previous words in \texttt{a0}, current words in
\texttt{a1}, result pointer in \texttt{a2}, and settings pointer in
\texttt{a3}. Settings contain two native unsigned doublewords, width and height,
each one through 64. This function also accepts smaller dimensions; MLV1 files remain 64 by 64.

Each input has exactly $wh$ native words and requires four-byte alignment. The
inputs may coincide. Compare all 32 bits, including a differing high bit; this
primitive does not validate serialized RGB colors. Settings and result require
eight-byte alignment. Their selected extents are sixteen and 32 bytes. Settings
and output are disjoint from each other and both input intervals. Do not add
partial-output alias support implicitly.

On success, return zero in \texttt{v0} and the number of unequal words in
\texttt{v1}. Publish four native doublewords $(x_0,y_0,x_1,y_1)$, where
$x_1,y_1$ are exclusive ends. Every changed position lies inside, and each
nonempty side is minimal. If inputs are equal, publish four zeros and count
zero. Input words and settings remain unchanged. Write no caller byte outside
the selected result.

Reject invalid dimensions or alignment with status one and count zero, without
reading image words or changing any caller byte. Test settings alignment before
loading settings. Preserve all saved registers, frame pointer, stack pointer and
the supplied return continuation; the incoming stack is sixteen-byte aligned.
Do no console or file I/O in the function.

The default demo has width two and height three. Indices one and four differ,
which means coordinates $(1,0)$ and $(0,2)$. The result is $(0,0,2,3)$ with
count two. Its rectangular area is six. Pause at
\texttt{bounds\_demo\_return} to inspect \texttt{v0/v1} and the native result
before the caller records status/count. A minimum/maximum linear-index interval
alone cannot determine the minimum horizontal bounds across multiple rows.

For one pass, keep a row and column cursor, changed count and four extrema.
Initialize minimum coordinates beyond any selected position, and maxima to zero;
handle the empty result explicitly. For another valid approach, first count and
locate changed rows, then find horizontal extrema in a second pass. The feedback
has an explicit finite execution budget that admits both approaches. It grades
exact coordinates and state rather than an instruction sequence or pass count.

\section{Check the extension and full codec separately}

Use \emph{Check this program} in the native lesson. The first stage retains
44 equality/alias cases, followed by 66 edge/row-crossing, 44 mixed/full-shape
and 28 boundary/reuse cases. The final assessment rechecks all 182 on this same
source in both byte orders. Edits make earlier execution evidence stale; follow
the renewal prompts. Saved predictions and hint history remain. This feedback
checks behavior and ownership without requiring a private label or exact loop.

\Needspace{10\baselineskip}
For standalone CLI work, run the independent feedback from the book repository:
\MARSCodeLabel{Check changed-pixel coordinates and ownership}
\begin{Verbatim}
java -jar "$CLI" grade \
  revision/labs/checks/video-change-bounds.kts \
  revision/labs/files/video-change-bounds.asm
\end{Verbatim}
It covers both byte orders, full-width invalid dimensions, word alignment
without extra vector alignment, coincident inputs, empty changes, single edges,
row crossings, the moving patch and repeated calls. Use sparse corners to show
that changed count and rectangle area can differ substantially. Repeat a
nonempty result followed by an unchanged pair to expose stale extrema.

The original codec uses its separate harness and domains:
\Needspace{8\baselineskip}
\MARSCodeLabel{Check one completed codec stage}
\begin{Verbatim}
python3 packages/034-lossless-video/grade.py \
  revision/labs/files/lossless-video.asm --domain seek \
  --cli "$CLI" --java "$JAVA"
\end{Verbatim}
Here \texttt{JAVA} is the full JDK 25 Java executable. Omit \texttt{--domain}
to check all eleven public domains. Reserve the \texttt{videoCheck\_} and
\texttt{videoPrivate\_} prefixes and their grading storage. The independent
bounds function has its own checks; the file-codec runner does not grade it.
Finally compare every decoded frame and every seek target with the generated
PPM previews, retaining the actual encoded file and source revision with your
observations.
